% ---------- Abstract ----------
\thispagestyle{plain}
\phantomsection\addcontentsline{toc}{chapter}{Abstract}
\begin{center}
  {\Large\bfseries ABSTRACT}
\end{center}

\vspace{0.8\baselineskip}

\begin{spacing}{1.3}
    \begin{otherlanguage}{english}
    \justifying
    \setlength{\parskip}{6pt plus 1pt minus 1pt}
    \setlength{\parindent}{0pt}

    Measuring fish from images would avoid removing them from the water, but photographing them through a tank introduces an error: light is refracted as it crosses water, glass and air, and the apparent size of the fish depends on its distance from the glass, which a single image cannot reveal. The objective was to compare, on the same anatomical points, four procedures for converting pixels to millimetres ---the scale of an ArUco marker attached to the glass (A), a homography from a marker coplanar with the replica (\Aprime), an underwater pinhole calibration (B) and a flat-port refractive model (C)--- and to quantify the length interval that results from an unknown depth. Six rigid replicas, measured with a calliper, were photographed at four distances from the glass, with three captures per condition (72 images). B and C require the depth to be known; with that information they obtained the lowest total-length errors (MAE 1.08 mm, 0.73\,\%, and 1.22 mm, 0.88\,\%), followed by \Aprime{} (3.29 mm) and A (37.43 mm); no difference between B and C was detected (sign test, $p = 1$). Without depth information, the interval of C spans 56 to 82\,\% of the length and contained the calliper length in all 72 captures; with depth bounded to $\pm 1$~cm it narrows to 4 to 5\,\%. It is concluded that, with known depth, the underwater calibration and the refractive model reduce the mean error below 1\,\%, and that without it depth is the dominant source of uncertainty.

    \textbf{Keywords:} \Keywords.
    \end{otherlanguage}
\end{spacing}
\clearpage
